\begin{figure}[tbp]
  \centering
  \includegraphics[width=1.0\linewidth,height=0.78\textheight,keepaspectratio]{figures/fig09.pdf}
  \caption{Illustrative examples of the effect of image features on candidate scores on the synthetic test sequence.
Reconstructed images are shown at \(-0.25\), \(0\), and \(+0.25\,\unit{ms}\) relative to the true collision time in (a), and at \(-0.50\), \(0\), and \(+0.50\,\unit{ms}\) relative to the candidate time in (b).
Panel (b) shows a shared-particle artifact, in which two reconstructed track IDs represent the same physical particle.
Red circles and blue squares mark actual observations of the terminating and continuing tracks, respectively; lines show their recent trajectories.
Each \(128\times128\)-pixel crop (\(1.28\times1.28\,\unit{mm}\)) is centered on the continuing track and reconstructed at its estimated depth; all six crops lie entirely within the recorded field of view.
For display, recorded pixels in each crop are linearly normalized between their first and 99th intensity percentiles; the scale bar is \ensuremath{500\,\unit{\mu m}}.
The diameters below (a) are ground-truth parent and daughter diameters.
(c) Candidate scores from the trajectory classifier and the image-feature ensemble, before relation-attention adjustment.
The true-event score changes from 0.999 to 0.998, whereas the artifact score decreases from 0.944 to 0.625.
The retained/rejected labels refer to the final candidate-selection policy.
These examples were selected for visibility after examining the results and illustrate individual candidates rather than an average effect over the test sequence.}
  \label{fig:image_feature_case_study}
\end{figure}
